\lysh{36777}{2}

\begin{alist}
\item Ισχύει ότι:
\begin{align*}
|x - 3| &\le 2 \Leftrightarrow -2 \le x - 3 \le 2 \Leftrightarrow -2 + 3 \le x - 3 + 3 \le 2 + 3 \Leftrightarrow 1 \le x \le 5. \\
|y - 6| &\le 4 \Leftrightarrow -4 \le y - 6 \le 4 \Leftrightarrow -4 + 6 \le y - 6 + 6 \le 4 + 6 \Leftrightarrow 2 \le y \le 10.
\end{align*}

\item Η περίμετρος ενός ορθογωνίου με διαστάσεις $2x$ και $y$ είναι $\Pi = 4x + 2y$.
Από το α) ερώτημα έχουμε:
$1 \le x \le 5 \Leftrightarrow 1 \cdot 4 \le 4x \le 5 \cdot 4 \Leftrightarrow 4 \le 4x \le 20$ (1)
$2 \le y \le 10 \Leftrightarrow 2 \cdot 2 \le 2y \le 10 \cdot 2 \Leftrightarrow 4 \le 2y \le 20$ (2)

Προσθέτοντας κατά μέλη τις ανισότητες (1) και (2) έχουμε:
\[4 + 4 \le 4x + 2y \le 20 + 20 \Leftrightarrow 8 \le \Pi \le 40.\]

Συνεπώς, η ελάχιστη τιμή που μπορεί να πάρει η περίμετρος είναι $8$, όταν $x = 1, y = 2$ και η μέγιστη τιμή που μπορεί να πάρει η περίμετρος είναι $40$, όταν $x = 5, y = 10$.
\end{alist}
